\documentclass[11pt,a4paper]{article}

\usepackage[margin=1in]{geometry}
\usepackage[T1]{fontenc}
\usepackage[utf8]{inputenc}
\usepackage{lmodern}
\usepackage{booktabs}
\usepackage{longtable}
\usepackage{array}
\usepackage{hyperref}
\usepackage{xcolor}

\hypersetup{
  colorlinks=true,
  linkcolor=blue,
  urlcolor=blue,
  pdftitle={POC Regression Report},
  pdfauthor={NLP Research Workspace}
}

\title{POC Regression Report\\Thai Medical OCR Pipeline Blocks}
\author{NLP\_RESEARCH\_TOPICS Workspace}
\date{2026-04-06}

\begin{document}
\maketitle

\section{Executive Summary}
This report summarizes the end-to-end regression run of six Proof-of-Concept (POC) modules used in the Thai medical OCR research pipeline. All modules are currently in a passing state after targeted bug fixes. The final state is:
\begin{itemize}
  \item 6/6 modules pass runnable tests.
  \item 4 modules required code fixes during regression.
  \item 1 module (PaddleOCR v5) required a backend migration from Paddle runtime to HuggingFace transformers.
\end{itemize}

\section{Environment}
\begin{itemize}
  \item OS: Linux
  \item GPU: NVIDIA GeForce RTX 5070 (Blackwell, SM\_120)
  \item CUDA target: cu128 (required for SM\_120 compatibility)
  \item Python/package manager: \texttt{uv}
  \item Workspace root: \texttt{/home/puem/Desktop/NLP\_RESEARCH\_TOPICS}
\end{itemize}

\section{Pipeline Context}
The tested modules map to the Thai medical OCR pipeline as follows:
\begin{itemize}
  \item Stage 2 (Layout Detection): DocLayout-YOLO
  \item Stage 3 (Text Recognition): EasyOCR Thai, PaddleOCR v5, GOT-OCR2
  \item Stage 4 (Post-correction): ByT5 OCR correction, WangchanBERTa correction
\end{itemize}

\section{Test Inputs}
Two main input classes were used:
\begin{enumerate}
  \item Synthetic samples generated by each module's \texttt{generate\_sample.py}
  \item A real academic paper image: \texttt{/home/puem/Desktop/NLP\_RESEARCH\_TOPICS/test\_image.png} (single-column page)
\end{enumerate}

\section{Final Regression Results}
\begin{longtable}{>{\raggedright\arraybackslash}p{2.8cm} >{\raggedright\arraybackslash}p{2.2cm} >{\raggedright\arraybackslash}p{2.2cm} >{\raggedright\arraybackslash}p{6.2cm}}
\toprule
Module & Stage & Status & Key Result \\
\midrule
\endfirsthead
\toprule
Module & Stage & Status & Key Result \\
\midrule
\endhead
DocLayout-YOLO & Stage 2 & PASS & Detected layout regions successfully on synthetic and real paper images. \\
EasyOCR Thai & Stage 3 & PASS & 42 regions (synthetic), 107 regions (real paper). JSON export stable after dtype fix. \\
PaddleOCR v5 (transformers backend) & Stage 3 & PASS & 8 regions (synthetic), 53 regions (real paper). Detection + recognition pipeline executes on GPU. \\
GOT-OCR2 & Stage 3 & PASS & 53 lines, 2978 characters on real paper in about 6.9 seconds. \\
ByT5 OCR correction & Stage 4 & PASS & 8 sample corrections executed; expected weak exact-match because model is not fine-tuned for this domain. \\
WangchanBERTa correction & Stage 4 & PASS & 8 sample mask-fills executed; 1/8 exact match on quick sanity set. \\
\bottomrule
\end{longtable}

\section{Issue Log and Fixes}
\subsection*{1) EasyOCR Thai JSON serialization failure}
\textbf{Symptom:} \texttt{TypeError: Object of type int32 is not JSON serializable} \\
\textbf{Root cause:} Bounding box values remained NumPy scalar types. \\
\textbf{Fix:} Cast bbox coordinates to Python floats before JSON write.

\subsection*{2) GOT-OCR2 API incompatibility}
\textbf{Symptom:} Missing \texttt{model.chat()} method in current transformers implementation. \\
\textbf{Root cause:} Runner assumed an older API style. \\
\textbf{Fix:} Switched to \texttt{AutoModelForImageTextToText} + \texttt{AutoProcessor} and generation-based inference.

\subsection*{3) Common CUDA device property bug}
\textbf{Symptom:} Access to \texttt{total\_mem} failed. \\
\textbf{Root cause:} Incorrect PyTorch CUDA property name. \\
\textbf{Fix:} Replaced with \texttt{total\_memory} in affected runners.

\subsection*{4) WangchanBERTa tokenizer dependency failure}
\textbf{Symptom:} Tokenizer initialization errors referencing \texttt{tiktoken}/SentencePiece extractor requirements. \\
\textbf{Root cause:} Missing protobuf runtime dependency in environment. \\
\textbf{Fix:} Added \texttt{protobuf} to project dependencies.

\subsection*{5) PaddleOCR v5 runtime failure and migration}
\textbf{Symptom:} Paddle runtime failed with low-level CPU path error (OneDNN / PIR conversion failure). \\
\textbf{Root cause:} Incompatibility in the tested Paddle stack for this environment. \\
\textbf{Fix:} Migrated module to HuggingFace transformers PP-OCRv5 models:
\begin{itemize}
  \item Detection model: \texttt{PaddlePaddle/PP-OCRv5\_server\_det\_safetensors}
  \item Recognition model: \texttt{PaddlePaddle/PP-OCRv5\_server\_rec\_safetensors}
\end{itemize}
Additionally repaired OpenCV wheel state and adapted detection parsing from polygon format \texttt{(N,4,2)} to axis-aligned crop boxes for recognition.

\section{Thai-Language Caveat}
The current PaddleOCR v5 transformers path uses the server recognition model with a non-Thai character inventory for the tested checkpoint. This is sufficient for English-heavy pages (as shown on the real paper test), but Thai-heavy medical forms may degrade in character fidelity. For production Thai documents, a Thai-specific recognition checkpoint (with compatible preprocessor metadata) is recommended.

\section{Reproducibility Commands}
From workspace root:
\begin{verbatim}
cd poc/doclayout-yolo && bash run.sh
cd ../easyocr-thai && bash run.sh
cd ../paddleocr-v5 && bash run.sh
cd ../got-ocr2 && bash run.sh
cd ../byt5-ocr-correction && bash run.sh
cd ../wangchanberta-correction && bash run.sh
\end{verbatim}

Real-image spot checks:
\begin{verbatim}
cd poc/paddleocr-v5
uv run python poc_runner.py --image ../../test_image.png --save-json
\end{verbatim}

\section{Conclusion}
The regression objective is met: all six modules now run successfully in the current environment. The highest-risk blocker (Paddle runtime instability) has been removed through a tested migration path, and the workspace is now in a stable state for next-stage integration and comparative evaluation experiments.

\end{document}
